\documentclass[11pt]{article}
\usepackage{series}
\usepackage{booktabs,tabularx,array}
\usepackage{xcolor}

\newcommand{\X}{\mathsf X}
\newcommand{\Y}{\mathsf Y}
\newcommand{\Z}{\mathsf Z}
\newcommand{\A}{\mathcal A}
\newcommand{\Rset}{\mathcal R}
\newcommand{\B}{\mathcal B}
\newcommand{\K}{\mathsf K}
\newcommand{\Ccap}{\mathcal C}
\newcommand{\Prob}{\mathbb P}
\newcommand{\E}{\mathbb E}
\newcommand{\R}{\mathbb R}
\newcommand{\N}{\mathbb N}
\newcommand{\MTT}{\mathrm{MTT}}

\title{\textbf{Projection-Limited Coherence}\\
\large Exact Criteria for Effective Dynamics and a Cross-Domain Research Program}
\author{Peter Nero}
\date{Version 2, July 2026}

\begin{document}
\maketitle

\begin{abstract}
Many sciences replace a detailed state by a smaller set of effective
variables.  This paper develops a substrate-neutral framework for asking
what such a projection does and does not imply.  For a measurable source
space $\X$, dynamics $\Phi$, observation map $\Pi:\X\to\Y$, admissible
domain $\A$, and declared margin $\Ccap$, we separate five questions:
whether deterministic dynamics descends to $\Y$, whether the descended map
is invertible, whether a projected stochastic process is Markov, whether
source information can be recovered, and whether the realizable image
$\Rset=\Pi(\A)$ obeys compatibility constraints.  Each question has its own
criterion.

We prove the exact fiber criterion for deterministic descent and the
corresponding lumpability criterion for stochastic kernels.  We show by
counterexample that a many-to-one projection, or even the absence of a
continuous section, need not make the effective dynamics irreversible.
Conversely, failure of the fiber criterion produces history dependence but
does not by itself prove computational irreducibility.  Information loss is
expressed by conditional entropy and data processing; monotone physical
entropy requires an additional dynamics and entropy-production law.
Likewise, a capacity is a control margin until a constitutive equation,
conservation law, or source term is supplied.

The resulting framework is useful precisely because it is selective.  It
organizes current Modal Triplet Theory constructions and provides a template
for models in physics, biology, cognition, artificial intelligence, and
collective systems.  The mathematical template transfers across domains;
physical conservation laws, entropy laws, selection principles, and claims
about consciousness or civilization do not.  Such applications are stated
as hypotheses with explicit validation obligations.
\end{abstract}

\section*{Version 2 revision note}

\begin{description}[style=nextline]
\item[Supersedes.]
Version 1, DOI \url{https://doi.org/10.5281/zenodo.18274610}.
\item[Reason.]
Version 1 treated many-to-one projection, finite capacity, and failure of a
global section as sufficient for irreversibility, entropy increase,
computational irreducibility, and cross-domain hidden relations.  Those
implications require different hypotheses and are false in that generality.
\item[Resolution.]
Version 2 replaces the universal chain by exact descent, lumpability,
recovery, image-geometry, information, and first-exit criteria.  Explicit
counterexamples mark the logical boundaries.  Capacity is kept separate
from conservation and entropy, and all nonphysical applications are
reclassified as model-building programs.
\item[Retained content.]
The central organizing idea survives: an effective description is governed
jointly by a source dynamics, a projection, a realizable image, and a
declared validity margin.  These data can generate loss of source
information, memory, compatibility constraints, and finite validity
windows, but only when the corresponding criterion is met.
\item[Open boundary.]
No universal domain-independent law selects $\Pi$, $\Ccap$, an entropy
functional, a probability measure, or a capacity dynamics.  Each physical
or empirical application must supply and test those objects.
\end{description}

\tableofcontents

\section{Why the distinctions matter}

An effective variable is not the same thing as a complete state.  A fluid
velocity field omits molecular coordinates; a detector record omits most of
the apparatus microstate; a summary statistic omits most of a data set.
This observation is elementary, but several different consequences are
often compressed into the phrase ``information was projected out.''

The compression is dangerous because the following statements are not
equivalent:
\begin{enumerate}
\item two source states have the same observed value;
\item the source dynamics induces an autonomous observed dynamics;
\item the induced observed dynamics is not invertible;
\item the observed stochastic process has memory;
\item the original source state cannot be recovered from the observation;
\item the observation has larger entropy at later times;
\item prediction has no computational shortcut.
\end{enumerate}
This paper gives each statement its own mathematical test.  The purpose is
not to deny that projection can contribute to irreversibility, memory, or
hidden constraints.  It is to identify the extra structure that turns that
possibility into a theorem.

The phrase \emph{projection-limited coherence} will mean the following
research pattern:
\[
  \text{source dynamics}
  \;+\;\text{observation map}
  \;+\;\text{realizable image}
  \;+\;\text{validity margin}.
\]
It is a template, not a substrate-independent equation of motion.

\section{Typed framework}

\subsection{Source, observation, and admissibility}

Let $(\X,\B(\X))$ and $(\Y,\B(\Y))$ be standard Borel spaces.  Let
\[
  \Phi:\X\longrightarrow\X
  \qquad\text{and}\qquad
  \Pi:\X\longrightarrow\Y
\]
be measurable.  The source dynamics may be discrete, or it may be the
time-$\Delta t$ map of a flow.  The observation map $\Pi$ need not be
linear or idempotent.  It may represent a quotient, a coarse-graining, a
measurement record, a feature map, or a projection followed by a
pushforward.

\begin{definition}[Admissible source and realizable image]
An \emph{admissible source domain} is a declared measurable subset
$\A\subseteq\X$.  Its realizable effective image is
\[
  \Rset_\Pi(\A):=\Pi(\A)\subseteq\Y.
\]
All effective constraints in this paper refer to this image or to a
probability measure supported on it.
\end{definition}

The distinction between $\Y$ and $\Rset_\Pi(\A)$ is essential.  Coordinates
on $\Y$ may look independent while the realizable image is a proper coupled
subset.  Equally, a many-to-one map can have
$\Rset_\Pi(\A)=\Y$, in which case no image constraint follows.

\subsection{Capacity as a declared margin}

\begin{definition}[Control margin]
A \emph{capacity} or \emph{control margin} is a function
\[
  \Ccap:\A\longrightarrow[0,\infty)
\]
whose operational meaning is fixed by a stated estimate.  Typical examples
are a spectral gap, distance to a failure set, inverse-condition margin,
coercivity constant, or certified truncation reserve.
\end{definition}

A nonnegative function called capacity is not automatically conserved,
extensive, additive, entropic, or physical.  Those properties are additional
claims.  Recent MTT capacity papers make this separation explicit by
distinguishing normalized margins, metric clearance, constitutive transport,
and first exit \cite{nero_capacity_invariant,nero_capacity_budget,
nero_capacity_dynamics}.

\subsection{Five maps that should not be conflated}

For later use, we name five distinct objects:
\begin{center}
\begingroup
\renewcommand{\arraystretch}{1.18}
\begin{tabularx}{\textwidth}{>{\raggedright\arraybackslash\bfseries}p{0.20\textwidth}>{\raggedright\arraybackslash}X}
\toprule
Object & Question answered \\
\midrule
$\Pi:\X\to\Y$ & What source information is retained? \\
$T:\Rset\to\Rset$ & Does source dynamics descend to an autonomous effective map? \\
$S:\Rset\to\X$ & Can one choose a representative source state? \\
$D:\Rset\to\X$ & Can one recover the actual source state on a declared class? \\
$\K(x,\cdot)$ & Does projected stochastic evolution depend only on the present effective state? \\
\bottomrule
\end{tabularx}
\endgroup
\end{center}
A section $S$ chooses one point in each fiber.  It does not recover which
point in that fiber actually occurred.  This simple distinction removes
several apparent paradoxes.

\section{When deterministic dynamics descends}

\subsection{The exact fiber criterion}

\begin{theorem}[Effective-dynamics fiber criterion]\label{thm:fiber}
Let $\Phi:\A\to\A$ and $\Pi:\A\to\Rset:=\Pi(\A)$ be maps.  There exists a
unique map $T:\Rset\to\Rset$ satisfying
\[
  T\circ\Pi=\Pi\circ\Phi
  \quad\text{on }\A
\]
if and only if
\[
  \Pi(x)=\Pi(x')
  \quad\Longrightarrow\quad
  \Pi(\Phi x)=\Pi(\Phi x')
  \qquad (x,x'\in\A).
\]
If the spaces and maps are measurable and the quotient is equipped with a
measurable structure for which the resulting $T$ is measurable, then this
is an autonomous measurable effective dynamics.
\end{theorem}

\begin{proof}
If $T$ exists, equal observed states have equal images after one step:
$\Pi(\Phi x)=T(\Pi x)=T(\Pi x')=\Pi(\Phi x')$.  Conversely, define
$T(y)=\Pi(\Phi x)$ for any $x$ with $\Pi(x)=y$.  The fiber condition makes
this independent of the chosen representative.  Uniqueness follows because
every $y\in\Rset$ has such a representative.
\end{proof}

This is the deterministic analogue of a sufficient statistic for evolution.
It says exactly when hidden coordinates can be forgotten without creating
state-dependent ambiguity.

\subsection{Three boundary examples}

\begin{example}[Many-to-one projection with invertible shadow]\label{ex:product}
Let $\X=\Y\times F$, let $\Pi(y,f)=y$, and let
\[
  \Phi(y,f)=(Ty,Rf)
\]
with $T$ and $R$ bijective.  Then $\Pi$ is many-to-one whenever $F$ has more
than one point, but the effective map is precisely the invertible map $T$.
Loss of source coordinates therefore does not imply irreversible effective
dynamics.
\end{example}

\begin{example}[No continuous section, yet invertible shadow]\label{ex:cover}
Let $\X=\Y=S^1$, $\Pi(z)=z^2$, and
$\Phi(z)=e^{i\alpha/2}z$.  Then
\[
  \Pi(\Phi z)=e^{i\alpha}\Pi(z),
\]
so the shadow map is the invertible rotation
$T(y)=e^{i\alpha}y$.  The double covering $\Pi$ has no continuous global
section, but this does not obstruct invertible shadow dynamics.  The
regularity class of a proposed section and the autonomy of the shadow are
different questions.
\end{example}

\begin{example}[Projection without autonomous descent]\label{ex:no-descent}
Let $\X=\{0,1\}^2$, let $\Pi(y,h)=y$, and let
$\Phi(y,h)=(h,y)$.  The source states $(0,0)$ and $(0,1)$ have the same
observation but evolve to observations $0$ and $1$.  The fiber criterion
fails, so no one-step map on $\Y$ represents the source dynamics.
\end{example}

Examples~\ref{ex:product}--\ref{ex:no-descent} are not pathologies.  They
show that projection supplies a question, while the source dynamics supplies
the answer.

\subsection{Sections, decoders, and actual recovery}

\begin{proposition}[Representative selection is not source recovery]
Suppose $S:\Rset\to\A$ is a right inverse of $\Pi$, so
$\Pi\circ S=\operatorname{Id}_{\Rset}$.  Then $S$ recovers every actual source state
$x\in\A$ from $\Pi(x)$ if and only if $\Pi$ is injective on $\A$ and
$S=\Pi^{-1}$ there.
\end{proposition}

\begin{proof}
Actual recovery requires $S(\Pi x)=x$ for every $x\in\A$, which makes $S$
a left inverse as well as a right inverse.  A map with a left inverse is
injective.  The converse is immediate.
\end{proof}

Thus failure of a section can matter for a chosen category or regularity
class, but existence of a section does not undo information loss inside a
fiber.  The appropriate recovery question must specify:
\begin{enumerate}
\item the source subset to be recovered;
\item the allowed decoder class;
\item exact or approximate recovery;
\item the norm, loss, or statistical risk;
\item whether side information is available.
\end{enumerate}
This is the same discipline used in statistics and information theory when
comparing experiments or sufficient representations
\cite{blackwell1953,cover_thomas}.

\section{Stochastic reduction and memory}

\subsection{The projected-kernel criterion}

Let $\K(x,\cdot)$ be a Markov kernel on $\X$.  Write
$\Pi^{-1}B=\{x:\Pi(x)\in B\}$ for $B\in\B(\Y)$.

\begin{theorem}[Projected Markov-kernel criterion]\label{thm:lump}
There exists a Markov kernel $\overline{\K}$ on $\Rset=\Pi(\X)$ such that
\[
  \overline{\K}(\Pi x,B)
  =\K(x,\Pi^{-1}B)
\]
for every source state $x$ and measurable $B\subseteq\Rset$ if and only if
\[
  \Pi(x)=\Pi(x')
  \quad\Longrightarrow\quad
  \K(x,\Pi^{-1}B)=\K(x',\Pi^{-1}B)
\]
for every measurable $B$.
\end{theorem}

\begin{proof}
Necessity follows because the right side must depend only on $\Pi x$.
Under the stated fiber constancy, the displayed formula defines
$\overline{\K}(y,B)$ independently of the representative $x$.
\end{proof}

For finite-state chains this is the familiar strong lumpability condition.
The general lesson is again exact: a projected process is Markov when the
transition probabilities of observable events are constant on observation
fibers.  Many-to-one projection neither guarantees nor forbids this
\cite{kemeny_snell,norris}.

\subsection{History dependence}

When Theorem~\ref{thm:lump} fails, a useful model may enlarge the state:
\[
  Y_n
  \quad\leadsto\quad
  (Y_n,Y_{n-1},\ldots,Y_{n-r+1})
\]
or introduce a predictive state inferred from the past.  A process has
finite Markov order $r$ only if its conditional law given the full past
equals the conditional law given the last $r$ observations.  Complete
connections and predictive-state constructions provide rigorous tools when
no small $r$ suffices \cite{walters1975,shalizi_crutchfield}.

\begin{remark}[What memory does not prove]
Failure of a one-step Markov representation on a chosen state space does
not prove that no finite augmentation exists, that simulation must follow
every microscopic step, or that prediction is computationally irreducible.
Those are complexity claims about a specified representation and algorithmic
task.
\end{remark}

\section{Realizable images and hidden compatibility}

Suppose an effective state has components
$\Y=\Y_1\times\cdots\times\Y_m$.  Define
\[
  \Rset=\Pi(\A),
  \qquad
  \Rset_j=\operatorname{pr}_j(\Rset).
\]

\begin{definition}[Support compatibility]
The effective variables obey a nontrivial \emph{support compatibility
constraint} when
\[
  \Rset\subsetneq\Rset_1\times\cdots\times\Rset_m.
\]
\end{definition}

\begin{proposition}[Recombination test]
The realizable image factors as
$\Rset=\prod_j\Rset_j$ if and only if every coordinatewise recombination of
values occurring in $\Rset$ also belongs to $\Rset$.
\end{proposition}

\begin{proof}
If $\Rset$ is the product of its coordinate projections, every
recombination lies in it.  Conversely, the recombination property places
every element of $\prod_j\Rset_j$ in $\Rset$.
\end{proof}

This proposition supplies a direct test.  Projection alone does not decide
its outcome.  Nor does support nonfactorization determine a probability
law.  If $\mu$ is a source measure and $\nu=\Pi_\#\mu$, statistical
independence requires
\[
  \nu=\nu_1\otimes\cdots\otimes\nu_m,
\]
which is a statement about the measure, not only its support.  Likewise, a
conditional-dependence graph is not automatically a physical interaction
graph.  These distinctions are developed in the companion image-geometry
paper \cite{nero_network_image}.

\section{Capacity, transport, and first exit}

\subsection{A margin becomes useful through an estimate}

Let $\Ccap:\A\to[0,\infty)$ and define the valid region
\[
  \A_+=\{x\in\A:\Ccap(x)>0\}.
\]
The first-exit time of a trajectory $x(t)$ is
\[
  \tau_\partial
  :=\inf\{t\geq0:\Ccap(x(t))=0\}.
\]
Without an evolution law for $\Ccap(x(t))$, this definition has no predicted
time scale.

\begin{proposition}[Elementary first-exit lower bound]\label{prop:exit}
Suppose $t\mapsto\Ccap(x(t))$ is absolutely continuous before first exit and
\[
  \left|\frac{d}{dt}\Ccap(x(t))\right|\leq L
\]
almost everywhere, with $\Ccap(x(0))=c_0>0$.  If $L>0$, then
\[
  \tau_\partial\geq\frac{c_0}{L}.
\]
If $L=0$, no finite exit occurs through this margin.
\end{proposition}

\begin{proof}
Absolute continuity gives
$\Ccap(x(t))\geq c_0-Lt$.  Reaching zero therefore requires
$t\geq c_0/L$.
\end{proof}

The proposition is modest but representative.  A capacity claim becomes
predictive only after a bound, balance law, transport equation, or
constitutive model connects the margin to dynamics.

\subsection{What is not automatic}

None of the following follows from the word ``capacity'':
\begin{itemize}
\item conservation of $\int\Ccap$;
\item a continuity equation;
\item monotone depletion;
\item additivity across subsystems;
\item equivalence to thermodynamic entropy;
\item a universal threshold shared by unrelated domains.
\end{itemize}
These may be imposed, derived, or tested in a particular model.  They cannot
be transferred from a physical theory to a biological or social model by
analogy alone.

\section{Information loss, entropy, and arrows of time}

\subsection{The exact information statement}

Let $X$ be a discrete source random variable and $Y=\Pi(X)$.  Then
\[
  H(X)=H(Y)+H(X\mid Y),
  \qquad
  H(Y)\leq H(X).
\]
The lost source information is the conditional entropy $H(X\mid Y)$.
For a stochastic channel $K$, the data-processing inequality gives
\[
  I(U;Y)\leq I(U;X)
\]
whenever $U\to X\to Y$ is a Markov chain
\cite{cover_thomas}.  Relative entropy also contracts under a Markov
kernel.

These statements compare information before and after a channel.  They do
not say that the marginal entropy of an observed system must increase with
physical time.  A deterministic coarse-graining of a discrete state
actually satisfies $H(\Pi(X))\leq H(X)$ at the instant of projection.
Thermodynamic entropy production requires additional ingredients such as a
dynamics, reference measure, coarse-grained equilibration, detailed-balance
structure, or fluctuation relation.

\subsection{Three notions of irreversibility}

\begin{definition}[Three irreversibility questions]
For a projected model, distinguish:
\begin{enumerate}
\item \emph{source loss}: $\Pi$ is noninjective on the relevant source set;
\item \emph{dynamical noninvertibility}: the descended map $T$ exists but is
not invertible;
\item \emph{predictive irreversibility}: retrodiction from available records
has a declared nonzero risk or entropy.
\end{enumerate}
\end{definition}

These notions can coincide in a model, but none is a synonym for the others.
Example~\ref{ex:product} has source loss and invertible effective dynamics.
A bijective observed map can still be hard to retrodict in the presence of
noise.  A stochastic process can be statistically time-asymmetric even when
its state transition matrix is invertible as a linear operator.

\subsection{Records and an arrow}

A record is a physical variable whose present value is statistically or
functionally coupled to a past event.  To derive an arrow of time from
records, a model must state:
\begin{enumerate}
\item the recording interaction;
\item the stable record algebra or state variables;
\item the preparation or boundary condition;
\item the direction in which record mutual information is generated;
\item the recovery or erasure protocol.
\end{enumerate}
MTT may supply a retarded branch orientation in a selected physical
construction, but projection alone does not select that orientation.

\section{Computational irreducibility needs a computational theorem}

A system is not computationally irreducible merely because its effective
process has memory.  A non-Markov process may have a compact predictive
state; a high-dimensional deterministic system may have an exactly soluble
observable; and a finite-state process can always be advanced by matrix
powers, although the cost may still be large.

\begin{definition}[Task-relative shortcut]
Fix an input encoding, an observable $f$, a time horizon $n$, an accuracy
$\varepsilon$, and a computational model.  A shortcut is an algorithm that
computes or approximates $f(\Phi^n x)$ with asymptotic cost smaller than the
declared step-by-step baseline.
\end{definition}

A no-shortcut theorem must therefore be a lower bound, reduction,
undecidability result, or complexity-theoretic statement for this typed
task.  Projection can motivate the search for such a theorem by exposing
hidden-state dependence.  It does not supply the theorem.

Predictive-state methods provide a constructive middle ground.  They ask
whether histories with the same conditional future distribution can be
merged into a smaller state.  The resulting causal states can reveal finite
or infinite predictive complexity without equating memory with
undecidability \cite{crutchfield_young,shalizi_crutchfield}.

\section{How the framework applies to MTT}

\subsection{A typed instantiation}

Within MTT, the abstract symbols can be assigned more structure:
\begin{center}
\begingroup
\renewcommand{\arraystretch}{1.12}
\begin{tabularx}{\textwidth}{>{\raggedright\arraybackslash\bfseries}p{0.20\textwidth}>{\raggedright\arraybackslash}X}
\toprule
Framework object & MTT role at the current declared tier \\
\midrule
$\X$ & An upper configuration, bundle, field, or finite operator domain
specified by the paper using it. \\
$\Pi$ & A selected coherent-sector projector, pushforward, truncation, or
observable map; these are not interchangeable without a commuting diagram. \\
$\A$ & The domain on which spectral, analytic, algebraic, or numerical
hypotheses have actually been verified. \\
$\Rset=\Pi(\A)$ & The set of effective data emitted by that selected map,
not all formally writable lower variables. \\
$\Ccap$ & A normalized margin, metric clearance, coercivity reserve,
spectral gap, or first-exit control quantity with its own estimate. \\
$T$ or $\overline{\K}$ & A descended deterministic or stochastic dynamics
only after the fiber or lumpability condition has been established. \\
\bottomrule
\end{tabularx}
\endgroup
\end{center}

This paper does not derive a universal MTT projector or a universal capacity
dynamics.  It provides the checklist by which a proposed MTT encoding can be
audited.  The current foundation supplies typed upper, coherent, and
effective layers \cite{nero_foundations}.  Current capacity work supplies
separate definitions for invariant margins and transport
\cite{nero_capacity_invariant,nero_capacity_budget,nero_capacity_dynamics}.
Current stochastic work distinguishes history dependence from quantum
algebra \cite{nero_indivisible}.

\subsection{Current claims and non-claims}

\begin{center}
\begingroup
\renewcommand{\arraystretch}{1.12}
\begin{tabularx}{\textwidth}{>{\raggedright\arraybackslash\bfseries}p{0.50\textwidth}>{\raggedright\arraybackslash}X}
\toprule
Statement & Status in this paper \\
\midrule
Fiber criterion for deterministic descent & Exact theorem. \\
Kernel criterion for Markov projection & Exact theorem. \\
Source information can be lost under a many-to-one map & Exact, relative
to a declared source set. \\
Projected dynamics is necessarily irreversible & False without further
hypotheses. \\
Projection necessarily increases physical entropy & False without an
entropy-production model. \\
Finite capacity obeys a universal conservation law & Not claimed. \\
MTT derives all effective physics from one universal projection & Not
claimed here. \\
Cross-domain systems instantiate the same physical law & Not claimed. \\
\bottomrule
\end{tabularx}
\endgroup
\end{center}

\section{Cross-domain use: analogy first, model second}

\subsection{What transfers}

The following abstract questions can be asked in any domain:
\begin{enumerate}
\item What is the source state?
\item What is observed or retained by $\Pi$?
\item What is the realizable image $\Rset$?
\item What estimate gives operational meaning to $\Ccap$?
\item Does the fiber or lumpability criterion hold?
\item Which prediction is made before examining the test data?
\end{enumerate}
This common syntax can be scientifically useful.  It can reveal that two
fields face structurally similar model-reduction problems.

\subsection{What does not transfer}

The following objects remain domain specific:
\begin{itemize}
\item conserved quantities and constitutive laws;
\item thermodynamic entropy and temperature;
\item biological fitness and metabolic cost;
\item cognitive reports and behavioral observables;
\item engineering failure criteria;
\item social institutions, incentives, and historical interventions.
\end{itemize}
Using the same word for two margins does not establish that they have the
same units, dynamics, or causal role.

\subsection{Domain-model ledger}

\begin{center}
\small
\begingroup
\renewcommand{\arraystretch}{1.10}
\begin{tabularx}{\textwidth}{>{\raggedright\arraybackslash\bfseries}p{0.13\textwidth}>{\raggedright\arraybackslash}p{0.20\textwidth}>{\raggedright\arraybackslash}p{0.22\textwidth}>{\raggedright\arraybackslash}X}
\toprule
Domain & Candidate source & Candidate projection & Required evidence \\
\midrule
Physics & A specified state, field, bundle, or operator model & A declared
measurement, quotient, or effective-field map & Commuting map, controlled
error, dynamics, units, and comparison with observables \\
Biology & Molecular, cellular, or population state & A phenotype, network,
or functional summary & Mechanistic model, perturbation data, uncertainty,
and held-out prediction \\
Cognition & Neural and bodily variables under an experimental protocol &
Reports or task-level latent variables & Operational definitions,
identifiability, intervention, and competing-model comparison \\
Artificial intelligence & Model state, training history, and environment &
Representations, outputs, or self-model variables & Reproducible benchmark,
ablation, distribution shift, and causal test \\
Collective systems & Agents, institutions, resources, and communication &
Macroscopic indicators or network summaries & Historical data,
counterfactual design, confound control, and domain-specific dynamics \\
\bottomrule
\end{tabularx}
\endgroup
\end{center}

\subsection{Consciousness and civilization}

Projection-limited coherence does not by itself define consciousness,
subjective time, agency, civilizational stability, or a solution to the
Fermi paradox.  Those topics can motivate candidate models, but each would
need operational variables and discriminating evidence.  In particular:
\begin{itemize}
\item a compressed self-model is not sufficient for consciousness;
\item memory in an observed process is not sufficient for agency;
\item a finite organizational margin is not a thermodynamic capacity;
\item failure of a social system is not a universal admissibility boundary;
\item absence of an observed extraterrestrial signal is not derived from
the abstract projection formalism.
\end{itemize}
The value of the framework here is methodological.  It says what would have
to be built before an analogy becomes a scientific model.

\section{Validation and falsification}

A projection-limited model should be accompanied by a seven-part
certificate:
\begin{enumerate}
\item \textbf{Source specification:} state space, units, and dynamics.
\item \textbf{Observation specification:} the map $\Pi$ and its uncertainty.
\item \textbf{Image test:} direct characterization of $\Rset=\Pi(\A)$ or a
certified approximation.
\item \textbf{Dynamics test:} fiber criterion, lumpability test, or an
explicit history-dependent alternative.
\item \textbf{Margin certificate:} definition of $\Ccap$ and the estimate it
controls.
\item \textbf{Prediction:} an observable not used to choose the model.
\item \textbf{Comparison:} uncertainty budget and performance against
plausible alternatives.
\end{enumerate}

The framework is falsified in an application when its declared map or margin
fails its own certificate, when an alleged image constraint admits observed
recombinations, when a claimed memory effect disappears under a sufficient
state augmentation, or when a domain-specific prediction fails.  The
abstract mathematics itself is not evidence for any particular empirical
instantiation.

\section{Discussion}

\subsection{A useful change of emphasis}

The strongest part of projection-limited coherence is not the claim that
projection explains everything.  It is the insistence that an effective
theory has an upstream source, a map, an image, and a validity domain.
Writing these objects explicitly often exposes where a proposed explanation
has skipped a step.

For example, if a projected process appears irreversible, one should first
ask whether an autonomous effective map exists.  If it does, test its
invertibility.  If it does not, model the hidden-state dependence rather
than calling the process fundamentally random.  If entropy is invoked,
identify the probability law and reference dynamics.  If capacity is
invoked, identify the norm or estimate it controls.  If two variables are
said to have a hidden relation, characterize the realizable image or the
joint measure.

\subsection{Why the corrected framework is stronger}

Removing universal implications does not make the framework empty.  It
makes it modular:
\[
\begin{array}{c}
\text{fiber test} \longrightarrow \text{autonomous deterministic shadow},\\
\text{lumpability test} \longrightarrow \text{autonomous Markov shadow},\\
\text{image test} \longrightarrow \text{compatibility constraints},\\
\text{information test} \longrightarrow \text{recoverability limits},\\
\text{margin estimate} \longrightarrow \text{controlled validity window}.
\end{array}
\]
Different applications may satisfy different rows.  The framework records
that fact without forcing all systems into one physical mechanism.

\subsection{Relationship to reduction and emergence}

The formalism is compatible with several established viewpoints.
Sufficient statistics and Blackwell comparison describe what information a
channel preserves \cite{blackwell1953}.  Lumpability describes when a
Markov chain descends to a partition \cite{kemeny_snell,norris}.
Information theory quantifies channel loss \cite{cover_thomas}.
Predictive-state methods ask which histories have the same future law
\cite{crutchfield_young,shalizi_crutchfield}.  Projection-limited coherence
does not replace these theories.  It provides a common typed ledger in which
their distinct obligations can be combined with MTT's source, admissibility,
and margin language.

\section{Conclusion}

Projection is not a universal generator of irreversibility, entropy,
complexity, or cross-domain law.  It is one component of an effective
description.  Exact consequences arise only after the source dynamics,
observation fibers, realizable image, probability law, and validity margin
are specified.

Version 2 establishes a corrected core:
\begin{enumerate}
\item deterministic descent is governed by fiber invariance;
\item stochastic descent is governed by kernel lumpability;
\item a section selects a representative but does not recover an unknown
source point in a nontrivial fiber;
\item compatibility is a property of the actual image or joint measure;
\item information loss does not by itself supply thermodynamic entropy
production;
\item memory does not by itself supply a no-shortcut theorem;
\item capacity is a control margin until a domain-specific law gives it
more structure.
\end{enumerate}

This is the defensible unifying statement.  MTT and other theories may
instantiate the template in multiple ways.  Biology, cognition, artificial
intelligence, and collective systems may also use it, but only through
independently specified and tested domain models.

\bibliographystyle{plain}
\bibliography{main}

\end{document}
